% The recursion behind the height bound, drawn as the objects it describes.
% T_h is built from T_{h-1} and T_{h-2} under one new root, so the node counts
% run 1, 2, 4, 7, 12 --- each one the sum of the previous two, plus the root.
%
% The macro doing the drawing is the same one the talk uses. It places each
% subtree by the golden ratio through Binet's formula, so the picture that
% argues for the golden ratio is itself laid out by it.
\begin{tikzpicture}[x = 1mm, y = 1mm]

  \setavlscale{3.1}{6.2}

  \foreach \h [count = \i from 0] in {0, 1, 2, 3, 4} {
    \begin{scope}[xshift = \i * 27mm]
      \resetavl
      \drawfib{\h}{0}{0}
      \node[font = \small] at (0, -33) {$T_{\h}$};
    \end{scope}
  }

  % node counts under each tree
  \foreach \n [count = \i from 0] in {1, 2, 4, 7, 12}
    \node[tlabel] at (\i * 27, -40) {$N = \n$};

  \node[tannot] at (54, -50)
    {each tree is the two before it, joined under one new root};

\end{tikzpicture}
